% Appendix C: theory behind the leakage-detection results (uses the theorem environments of main.tex).

\section{Theory of the leakage-detection results}\label{app:theory}

This appendix derives the three results used in Sections~\ref{sec:leakresults} and~\ref{sec:gradedleak}: the optimism of pooled multi-draw planning, detection from a committed score ledger, and the sensitivity of held-out error to selection leakage.

\subsection{Replacement testing over a population of entries}

Fix a replacement strategy. For a task-partition entry \(w\) (a task, its five-role partition, its reference payload, and its frozen contexts), let \(\pi_w\in[0,1]\) be the per-draw probability that a replacement rerun alters the payload \(P\), and let \(\Pi\) be the distribution of \(\pi_w\) over a population of entries, with expectation \(\mathbb E_\Pi[\cdot]\).

\begin{theorem}[Population detection under heterogeneous sensitivity]\label{thm:mixture}
Let \(D(t)=\mathbb E_{\Pi}\big[1-(1-\pi_w)^t\big]\) be the population-average probability that \(t\ge1\) independent replacement draws detect a violation, and let \(\bar\pi=\mathbb E_\Pi[\pi_w]\). Then, for every integer \(t\ge1\),
\begin{equation}\label{eq:jensengap}
D(t)\le 1-(1-\bar\pi)^t,
\end{equation}
with equality for \(t\ge2\) if and only if \(\pi_w=\bar\pi\) holds \(\Pi\)-almost surely.
\end{theorem}

\begin{pf}
Equality is trivial for \(t=1\). For \(t\ge2\), \(\phi_t(u)=1-(1-u)^t\) is strictly concave on \([0,1]\) because \(\phi_t''(u)=-t(t-1)(1-u)^{t-2}<0\) on \([0,1)\); Jensen's inequality gives \(D(t)=\mathbb E_\Pi[\phi_t(\pi_w)]\le\phi_t(\bar\pi)\), with equality if and only if \(\pi_w\) is \(\Pi\)-almost surely constant.
\end{pf}

Theorem~\ref{thm:mixture} explains the optimism of population-pooled multi-draw planning when \(\pi_w\) varies across entries. Entry-wise estimation avoids this aggregation bias: if \(M\ge t\) independent draws on entry \(w\) yield \(J\) payload changes, \(\binom{M-J}{t}/\binom{M}{t}\) is an unbiased estimator of the miss probability \((1-\pi_w)^t\), because it averages the no-change indicators over all size-\(t\) subsets of the \(M\) draws. With \(M=t=3\), as in Figure~\ref{fig:leakage}(B), it reduces to the indicator that none of the three draws changed the payload.

\subsection{Score-ledger detection and predictive sensitivity}

\begin{theorem}[Detection from committed ledger collision mass]\label{thm:ledger}
Fix the construction observations, contexts, reference calibration data \(D_{\rm cal}\), and leakage weight \(\lambda\). Let \(\mathcal L_\lambda(D)\) be the canonical score-ledger byte string produced using calibration data \(D\), including the binary64 keys of Eq.~\eqref{eq:gradedkey}, so that \(\mathcal L_\lambda(D_{\rm cal})\) is the ledger committed in the reference payload \(P\). For a replacement draw define the collision mass
\begin{equation}\label{eq:ledgercollision}
c_w=\Pr\big(\mathcal L_\lambda(\widetilde D_{\rm cal})=\mathcal L_\lambda(D_{\rm cal})\big).
\end{equation}
The ledger-comparison detector has single-draw detection probability \(1-c_w\), the complete-payload detector has \(\pi_w\ge1-c_w\), and \(t\) independent replacement reruns satisfy
\begin{equation}\label{eq:ledgermiss}
\Pr\big(P_1=\cdots=P_t=P\big)\le c_w^t.
\end{equation}
\end{theorem}

\begin{pf}
Ledger comparison changes exactly on the complement of the event in Eq.~\eqref{eq:ledgercollision}. Since the complete ledger is serialised in \(P\), a changed ledger implies a changed payload. Thus a payload miss implies a ledger collision. Independence of the replacement draws makes the probability of \(t\) ledger collisions equal to \(c_w^t\), proving Eq.~\eqref{eq:ledgermiss}.
\end{pf}

For \(n=|D_{\rm cal}|\), bootstrap draws and response permutations give the exact finite-distribution expressions
\begin{align}
c_w^{\rm boot}
&=n^{-n}\sum_{(i_1,\ldots,i_n)\in\{1,\ldots,n\}^n}
\mathbf1\{\mathcal L_\lambda(D_{i_1,\ldots,i_n})=\mathcal L_\lambda(D_{\rm cal})\},\label{eq:bootcollision}\\
c_w^{\rm perm}
&=(n!)^{-1}\sum_{\tau\in\mathfrak S_n}
\mathbf1\{\mathcal L_\lambda(X,y_\tau)=\mathcal L_\lambda(D_{\rm cal})\},\label{eq:permcollision}
\end{align}
where \(D_{i_1,\ldots,i_n}\) is the ordered bootstrap sample, \((X,y_\tau)\) denotes the calibration covariates with responses permuted by \(\tau\), and \(\mathfrak S_n\) is the permutation group. Gaussian response noise uses the same collision-mass definition with its noise distribution. Paired-row shuffling tests the dependence of binary64 score reductions on row order. The complete ledger detected every leaking rerun reported in Section~\ref{sec:gradedleak}.

\begin{theorem}[Sensitivity asymmetry between structural and predictive instruments]\label{thm:asym}
\ifdefined\ReviewVersion\mbox{}\par\fi
Let \(m_0=\arg\min_{m\in\mathcal M_C} A(m)\) be the compliant winner and \(m_\lambda=\arg\min_{m\in\mathcal M_C}\widehat K_\lambda(m)\) the winner under the recorded key \(\widehat K_\lambda\), for \(\lambda\in[0,1)\). Write \(\eta_\lambda=\sup_{m\in\mathcal M_C}|\widehat K_\lambda(m)-K_\lambda(m)|\) for the rounding error relative to Eq.~\eqref{eq:gradedkey}. Let \(W=\max_{m\in\mathcal M_C} B(m;D_{\rm cal})-\min_{m\in\mathcal M_C} B(m;D_{\rm cal})\) be the range of the calibration score over the bank, let \(R(m)\) be the population RMSE of candidate \(m\), and let \(\omega=\sup_{m\in\mathcal M_C}|A(m)-R(m)|\). Then
\begin{align}
A(m_\lambda)-A(m_0)&\le\frac{\lambda\{B(m_0;D_{\rm cal})-B(m_\lambda;D_{\rm cal})\}+2\eta_\lambda}{1-\lambda}\notag\\
&\le\frac{\lambda W+2\eta_\lambda}{1-\lambda},\label{eq:asymA}\\
R(m_\lambda)-R(m_0)&\le\frac{\lambda W+2\eta_\lambda}{1-\lambda}+2\omega,\label{eq:asymupper}\\
R(m_\lambda)-R(m_0)&\ge-2\omega.\label{eq:asymlower}
\end{align}
\end{theorem}

\begin{pf}
Optimality for the recorded key gives \(K_\lambda(m_\lambda)\le\widehat K_\lambda(m_\lambda)+\eta_\lambda\le\widehat K_\lambda(m_0)+\eta_\lambda\le K_\lambda(m_0)+2\eta_\lambda\). Substitution of Eq.~\eqref{eq:gradedkey} and division by \(1-\lambda\) yield Eq.~\eqref{eq:asymA}. Combining Eq.~\eqref{eq:asymA} with \(R(m_\lambda)\le A(m_\lambda)+\omega\) and \(-R(m_0)\le -A(m_0)+\omega\) yields Eq.~\eqref{eq:asymupper}, and combining \(A(m_0)\le A(m_\lambda)\) with the same deviation bound yields Eq.~\eqref{eq:asymlower}.
\end{pf}

Theorem~\ref{thm:asym} bounds the candidate-level population-risk shift by \((\lambda W+2\eta_\lambda)/(1-\lambda)+2\omega\). Empirical test RMSE also varies with the test sample. Defining \(m_1=\arg\min_m B(m;D_{\rm cal})\) and \(\omega_B=\sup_{m\in\mathcal M_C}|B(m;D_{\rm cal})-R(m)|\) gives \(R(m_1)-\min_m R(m)\le2\omega_B\). By Theorem~\ref{thm:ledger}, ledger detection instead depends directly on the collision mass of the committed construction scores.
